\documentclass[12pt]{article}
\usepackage[utf8]{inputenc}
\usepackage[T1]{fontenc}
\usepackage{hyperref}
\title{T\`ay Sa Pa Language Resources for K-12 Curriculum Development}
\author{Compiled by Claude Sonnet 5 (claude-sonnet-5) for Andr\'as Kornai}
\date{}
\begin{document}
\maketitle

\section{Introduction}
T\`ay Sa Pa (also Sapa or T\`ay Muang Bo; ISO 639-3 tys, Glottocode tays1238) is a small Southwestern Tai language spoken in Sa Pa, L\`ao Cai Province, northern Vietnam. Wikipedia gives about 300 speakers (2002). Joshua Project lists 400 people. Vietnamese authorities classify the speakers as T\`ay. The searches found only descriptive and reference sources.

\section{Primary Educational Resources}
No primer, textbook or school programme turned up. Joshua Project describes the language as unwritten.

\section{Linguistic and Reference Materials}
\begin{itemize}
\item Sapa language (Wikipedia): \url{https://en.wikipedia.org/wiki/Sapa_language}
\item Journal article on the phonology of the T\`ay Muang Bo lect (Sa Pa, L\`ao Cai), by S\^am C\^ong Danh. It describes the work as preliminary: \url{https://vjol.info.vn/dhcnhn/article/view/109852}
\end{itemize}

\section{Religious and Community Resources}
\begin{itemize}
\item Joshua Project profile of the Tay Sa Pa people. It reports that Bible translation has started: \url{https://joshuaproject.net/people_groups/19578}
\item Joshua Project language page for tys: \url{https://joshuaproject.net/languages/tys}
\end{itemize}

\section{Digital and Audio Resources}
\begin{itemize}
\item Global Recordings Network page for tys. I did not open it to check what recordings it holds: \url{https://glorecnet-website.k.globalrecordings.net/en/language/tys}
\end{itemize}

\section{Curriculum Development Context}
The language has few speakers, no established writing system according to the sources found, and little documentation beyond the preliminary phonological study. Any curriculum would first need documentation and orthography work with the community. The Vietnamese alphabet is the likely basis, as Wikipedia lists it.
\end{document}